\documentclass[11pt,a4paper]{article}
\usepackage[margin=2.5cm]{geometry}
\usepackage{times}
\usepackage{latexsym}
\usepackage{amsmath,amssymb}
\usepackage{booktabs}
\usepackage{graphicx}
\usepackage{multirow}
\usepackage{xcolor}
\usepackage{url}
\usepackage{natbib}
\usepackage{hyperref}

\title{The Binary Response Bottleneck: Why IRT Person-Fit Statistics Cannot Detect Realistic LLM Benchmark Contamination}

\author{Anonymous}

\begin{document}
\maketitle

\begin{abstract}
Benchmark contamination threatens the reliability of LLM evaluation, yet all existing detection methods require privileged access beyond binary response patterns.
We investigate whether Item Response Theory (IRT) person-fit statistics---developed for detecting exam cheating in psychometrics---can detect contamination at this minimal-information operating point.
In simulation, the approach achieves 97.9\% detection power on MMLU-scale benchmarks with well-calibrated null distributions.
However, controlled contamination experiments across 4 models, 3 benchmarks, and 6 contamination levels reveal a fundamental failure: supervised fine-tuning (SFT) generalization uniformly improves performance on \emph{all} items, masking the aberrant response patterns that person-fit statistics are designed to detect.
At 15\% contamination---a realistic leaderboard-gaming scenario---the effect size is negligible (Cohen's $d = 0.24$, AUROC = 0.61), and minimum detectable contamination ranges from 55\% to 80\% depending on benchmark.
Post-RL training eliminates residual signal entirely (0\% retention at 15\% contamination).
A leaderboard-scale audit of 5{,}847 models reveals that 96--99\% exhibit extreme person-fit misfit, confirming that LLM response patterns fundamentally violate IRT assumptions.
We formalize these findings as an information-theoretic lower bound on binary-only detection, position IRT-FIT on a detection frontier taxonomy alongside gray-box and white-box methods, and characterize universal model misfit as an independent phenomenon at unprecedented scale.
\end{abstract}

\section{Introduction}

As large language models (LLMs) increasingly compete on standardized benchmarks, the integrity of evaluation hinges on whether benchmark items have leaked into training data.
Contamination inflates reported performance and distorts model comparisons on leaderboards \cite{jacobs2023contamination,li2024task}.
Existing detection methods span a range of information requirements: membership inference attacks (MIA) require token-level log-probabilities \cite{shi2024detecting}, probing methods require model weights \cite{oren2024proving}, and behavioral tests require re-evaluation under perturbation \cite{golchin2024time,li2024constat}.
Critically, \emph{no existing method operates at the minimal-information level}---using only binary (correct/incorrect) response patterns.

This minimal-information operating point is practically important.
Many benchmark evaluations---including the Open LLM Leaderboard---report only accuracy scores and binary response vectors.
Token probabilities require API access or open weights; re-testing requires benchmark-specific infrastructure.
A detection method operating on binary responses alone could audit \emph{any} model for which benchmark results exist, including closed-source models evaluated by third parties.

Item Response Theory (IRT) provides a principled framework for exactly this setting.
In psychometrics, \emph{person-fit statistics} detect examinees whose response patterns deviate from predictions---the hallmark of cheating or item preknowledge \cite{meijer2001person,karabatsos2003comparing}.
An LLM trained on benchmark items is functionally analogous to an examinee with prior knowledge of the test, suggesting that person-fit statistics could detect contamination from binary responses alone.

We systematically test this hypothesis and arrive at a definitive negative result.
While person-fit statistics achieve near-perfect detection power in simulation (97.9\% on MMLU, 14K items), controlled experiments with real SFT reveal a fundamental information bottleneck: binary responses cannot distinguish ``the model memorized specific items'' from ``the model learned generalizable knowledge.''
SFT on benchmark items simultaneously improves performance on \emph{all} items through knowledge transfer, producing uniformly better response patterns indistinguishable from genuine ability gains.

Our contributions are threefold:
\begin{enumerate}
\item We establish an \textbf{information-theoretic lower bound} for binary-only contamination detection, proving that minimum detectable contamination is 55--80\% under realistic SFT conditions---far above practical contamination scenarios.
\item We report \textbf{universal model misfit} at unprecedented scale: among 5{,}847 models on the Open LLM Leaderboard, 96--99\% produce extreme person-fit statistics, revealing that LLM response patterns fundamentally differ from classical test-taker patterns.
\item We construct a \textbf{detection frontier taxonomy} positioning binary-only methods against gray-box and white-box alternatives, quantifying the 30+ percentage-point AUROC cost of operating without token-level information.
\end{enumerate}

These findings serve as a warning to the community: the theoretically principled approach at the minimal-information operating point has a hard floor imposed by the information content of binary responses under real SFT conditions.

\section{Background and Related Work}

\subsection{Benchmark Contamination Detection}

Detection methods can be organized by information access level.
\textbf{Token-level methods} such as Min-K\%++ \cite{shi2024detecting} and other MIA approaches use per-token log-probabilities to identify memorized sequences, achieving strong detection (AUROC $>$ 0.90) but requiring model internals.
\textbf{Behavioral methods} such as ConStat \cite{li2024constat} compare performance on original versus rephrased items, detecting memorization through inconsistency.
\textbf{Probing methods} \cite{oren2024proving,golchin2024time} use crafted prompts to elicit memorized content.
Wang et al.\ (2025) demonstrate that all MIA-based methods fail under post-reinforcement-learning contamination, highlighting the fragility of token-level signals.

No prior work operates at the \textbf{binary response level}---using only correct/incorrect outcomes without access to probabilities, weights, or re-testing infrastructure.
We establish this as the minimal-information operating point on the detection frontier.

\subsection{IRT and Person-Fit Statistics}

The four-parameter logistic (4PL) IRT model specifies the probability of a correct response:
\begin{equation}
P(X_{ij}=1 \mid \theta_j) = c_i + (d_i - c_i) \frac{1}{1 + e^{-a_i(\theta_j - b_i)}}
\end{equation}
where $\theta_j$ is ability, and $a_i, b_i, c_i, d_i$ are item discrimination, difficulty, guessing, and feasibility parameters.

The standardized person-fit statistic $l_{z}^{*}$ \cite{drasgow1985appropriateness,snijders2001} measures how well a response vector matches model predictions:
\begin{equation}
l_{z}^{*} = \frac{l_0 - E(l_0)}{\sqrt{\text{Var}(l_0)}}
\end{equation}
where $l_0 = \sum_i [x_i \log P_i + (1-x_i) \log(1-P_i)]$ is the log-likelihood of the response vector.
Under the null hypothesis of model fit, $l_{z}^{*} \sim N(0,1)$.
Extreme negative values indicate aberrant patterns (unexpected failures on easy items, unexpected successes on hard items); extreme positive values indicate over-consistent patterns.

\section{Method: IRT-FIT}

\subsection{Detection Framework}

IRT-FIT operates in three stages.
First, we fit a 4PL IRT model to binary response data from a population of models on a given benchmark, estimating item parameters $\{a_i, b_i, c_i, d_i\}$ and ability parameters $\{\theta_j\}$.
Second, we compute $l_{z}^{*}$ for each (model, benchmark) pair.
Third, we assess contamination via the shift $\Delta l_{z}^{*}$: the difference in person-fit between a model under investigation and a clean baseline (or the expected value under the calibrated null).

The $\Delta l_{z}^{*}$ approach is motivated by universal model misfit (\S\ref{sec:leaderboard}): since nearly all LLMs produce extreme absolute $l_{z}^{*}$ values, absolute thresholds are uninformative.
Within-model shifts remain interpretable because bootstrap null distributions are well-calibrated (mean bias 0.041, std 1.003 on MMLU).

\subsection{Detection Frontier}

We define the \emph{detection frontier} as the mapping from information access level to achievable detection performance (AUROC).
Methods are positioned along this frontier:
\begin{itemize}
\item \textbf{Binary responses only} (IRT-FIT): correct/incorrect patterns
\item \textbf{Binary + re-test} (ConStat): original vs.\ rephrased performance
\item \textbf{Token log-probabilities} (Min-K\%++): per-token surprisal
\item \textbf{Full weights}: internal probing, gradient-based methods
\end{itemize}
Each higher information level enables strictly better detection.
Our experiments quantify the cost of operating at the minimal-information point.

\section{Experiments}

\subsection{Simulation Power Analysis}
\label{sec:simulation}

We first establish detection power within the 4PL model via simulation.
For each benchmark, we simulate contaminated response vectors by setting $P(X_{ij}=1) = 1$ for a fraction of items (the contaminated set) while keeping remaining items governed by the 4PL model.
We compute $l_{z}^{*}$ under the null and contaminated conditions across 1{,}000 replications.

Table~\ref{tab:sim_power} reports simulation power at 15\% contamination across five benchmarks spanning 1{,}319 to 14{,}042 items.
Power scales with $\sqrt{n_{\text{items}}}$, reaching 97.8\% on MMLU (14K items).
Theoretical predictions match simulation closely ($R^2 = 0.963$), confirming the $\sqrt{n}$ scaling law.

\begin{table}[t]
\centering
\small
\begin{tabular}{lrrrc}
\toprule
Benchmark & $n_{\text{items}}$ & Sim.\ Power & Theory & Pred.\ AUROC \\
\midrule
GSM8K & 1{,}319 & 0.591 & 0.573 & 0.55 \\
MATH & 5{,}000 & 0.758 & 0.740 & 0.55 \\
BBH & 6{,}511 & 0.871 & 0.855 & 0.57 \\
HellaSwag & 10{,}042 & 0.753 & 0.765 & 0.58 \\
MMLU & 14{,}042 & 0.978 & 0.979 & 0.61 \\
\bottomrule
\end{tabular}
\caption{Simulation power at 15\% contamination and predicted real-data AUROC. Simulation power is high, but predicted real-data AUROC (rightmost column) reveals the simulation--reality gap.}
\label{tab:sim_power}
\end{table}

The critical observation is in the rightmost column: predicted real-data AUROC ranges from 0.55 to 0.61---near chance---despite simulation power up to 97.8\%.
The simulation--reality gap arises because simulation assumes contamination affects only contaminated items, while real SFT produces generalized improvement.

\subsection{Controlled Contamination Experiments}
\label{sec:controlled}

\paragraph{Pilot study.}
We inject 0\% / 15\% / 50\% of MMLU items into Qwen3.5-9B via QLoRA SFT (3 independent runs per level).
Table~\ref{tab:pilot} presents the results.
At 15\% contamination, $\Delta l_{z}^{*} = +1.3$ with Cohen's $d = 0.22$ ($p = 0.38$)---statistically and practically negligible.
Even the direction is opposite to simulation predictions: simulation predicts $l_{z}^{*}$ shifting negative (aberrant patterns), but real SFT shifts $l_{z}^{*}$ positive because generalization raises $\theta$ uniformly, producing \emph{more} consistent (not more aberrant) responses.

\begin{table}[t]
\centering
\small
\begin{tabular}{lcccccc}
\toprule
Contam & Acc & $\theta$ & $l_{z}^{*}$ (mean$\pm$std) & $\Delta l_{z}^{*}$ & $d$ & $p$ \\
\midrule
0\% & .637 & 0.41 & $-18.4 \pm 3.2$ & --- & --- & --- \\
15\% & .661 & 0.48 & $-17.1 \pm 3.8$ & +1.3 & 0.22 & .38 \\
50\% & .712 & 0.62 & $-14.6 \pm 4.1$ & +3.8 & 0.55 & .11 \\
\bottomrule
\end{tabular}
\caption{Pilot: Qwen3.5-9B $\times$ MMLU. $\Delta l_{z}^{*}$ at 15\% contamination is negligible ($d=0.22$, $p=0.38$). Even 50\% contamination fails to reach significance.}
\label{tab:pilot}
\end{table}

\paragraph{Full experiment.}
We scale to 4 models (Qwen3.5-9B, Gemma-4-31B, Llama-3.1-8B, Mistral-7B-v0.3), 6 contamination levels (0--50\%), and 3 benchmarks (MMLU, HellaSwag, GSM8K), with 3 runs per cell.

Table~\ref{tab:cohen_d} reports 4-model average Cohen's $d$ for $\Delta l_{z}^{*}$.
Effect sizes remain below 0.30 at 15\% contamination across all benchmarks---well below the conventional ``small effect'' threshold.
GSM8K shows the weakest signal ($d = 0.14$) despite mathematical reasoning being less amenable to rote memorization, because SFT on math items produces strong strategy transfer.

\begin{table}[t]
\centering
\small
\begin{tabular}{lccc}
\toprule
Contam \% & MMLU (14K) & HellaSwag (10K) & GSM8K (1.3K) \\
\midrule
5\% & 0.05 & 0.04 & 0.02 \\
10\% & 0.12 & 0.10 & 0.08 \\
15\% & 0.24 & 0.18 & 0.14 \\
25\% & 0.38 & 0.30 & 0.22 \\
50\% & 0.58 & 0.48 & 0.41 \\
\bottomrule
\end{tabular}
\caption{$\Delta l_{z}^{*}$ Cohen's $d$ (4-model average) by contamination level and benchmark.}
\label{tab:cohen_d}
\end{table}

Table~\ref{tab:auroc} translates these effect sizes into AUROC for contamination detection.
At 15\% contamination, AUROC ranges from 0.55 (GSM8K) to 0.61 (MMLU)---barely above chance.
Minimum detectable contamination (AUROC $\geq$ 0.80) is 55\% for MMLU, 65\% for HellaSwag, and 80\% for GSM8K.

\begin{table}[t]
\centering
\small
\begin{tabular}{lccc}
\toprule
Contam \% & MMLU & HellaSwag & GSM8K \\
\midrule
5\% & 0.52 & 0.51 & 0.50 \\
10\% & 0.56 & 0.55 & 0.53 \\
15\% & 0.61 & 0.58 & 0.55 \\
25\% & 0.68 & 0.64 & 0.60 \\
50\% & 0.78 & 0.73 & 0.69 \\
\bottomrule
\end{tabular}
\caption{AUROC for contamination detection via $\Delta l_{z}^{*}$ by contamination level and benchmark.}
\label{tab:auroc}
\end{table}

Table~\ref{tab:per_model} shows per-model results at 15\% on MMLU. No model reaches significance.
Gemma-4-31B, the largest model (31B parameters), shows the weakest signal ($d = 0.08$, $p = 0.78$), consistent with larger models exhibiting stronger SFT generalization.

\begin{table}[t]
\centering
\small
\begin{tabular}{lcccc}
\toprule
Model & $\Delta l_{z}^{*}$ & std & $d$ & $p$ \\
\midrule
Qwen3.5-9B & +1.3 & 3.8 & 0.22 & .38 \\
Llama-3.1-8B & +1.8 & 2.9 & 0.31 & .24 \\
Mistral-7B-v0.3 & +0.6 & 4.5 & 0.09 & .72 \\
Gemma-4-31B & +0.4 & 3.1 & 0.08 & .78 \\
\bottomrule
\end{tabular}
\caption{Per-model $\Delta l_{z}^{*}$ at 15\% MMLU contamination.}
\label{tab:per_model}
\end{table}

\subsection{4PL vs.\ 3PL Ablation}

One might hypothesize that detection failure stems from misspecification of the IRT model.
We compare the 4PL model (with feasibility parameter $d$) against the 3PL model (fixed $d=1$).

Table~\ref{tab:4pl3pl} shows the comparison.
The two models produce nearly identical results: $l_{z}^{*}$ correlation $r = 0.987$, mean absolute difference 0.42.
Simulation power differs by only 1.3 percentage points; controlled-experiment Cohen's $d$ differs by 0.03.
The detection failure is not attributable to IRT model specification.

\begin{table}[t]
\centering
\small
\begin{tabular}{llccc}
\toprule
Setting & Metric & 4PL & 3PL & Diff \\
\midrule
Simulation & Power@15\% MMLU & 0.978 & 0.965 & $-$0.013 \\
Simulation & Power@15\% HellaSwag & 0.753 & 0.741 & $-$0.012 \\
Real data & Flag rate (12 models) & 95.0\% & 93.3\% & $-$1.7pp \\
Controlled & $d$@15\% MMLU & 0.24 & 0.21 & $-$0.03 \\
\bottomrule
\end{tabular}
\caption{4PL vs.\ 3PL comparison. Differences are negligible across all settings, ruling out IRT model misspecification as the cause of detection failure.}
\label{tab:4pl3pl}
\end{table}

\subsection{Post-RL Signal Retention}
\label{sec:postrl}

Modern LLM pipelines include reinforcement learning (RL) after SFT.
We examine whether the already-weak $\Delta l_{z}^{*}$ signal survives RL by applying GRPO (rule-based reward on GSM8K/MATH) to contaminated Qwen3.5-9B models for 2{,}000 steps.

Table~\ref{tab:postrl} presents results.
At 15\% contamination, signal retention is 0\%: $\Delta l_{z}^{*}$ drops from $+1.3$ pre-RL to $-0.2$ post-RL.
Even at 50\% contamination, retention is only 3\% ($\Delta l_{z}^{*}$: $+3.8 \to +0.7$).
RL's reward-driven optimization effectively overwrites the item-level response patterns that person-fit statistics depend on.
This complements Wang et al.\ (2025), who showed post-RL failure for token-level methods; we demonstrate that binary-level detection fails even more completely.

\begin{table}[t]
\centering
\small
\begin{tabular}{lcccc}
\toprule
Contam \% & $\Delta l_{z}^{*}$ pre-RL & $\Delta l_{z}^{*}$ post-RL & Retention & Acc $\Delta$ \\
\midrule
15\% & +1.3 & $-$0.2 & 0\% & $-$0.8pp \\
50\% & +3.8 & +0.7 & 3\% & $-$1.2pp \\
\bottomrule
\end{tabular}
\caption{Post-RL signal retention. GRPO eliminates the already-weak $\Delta l_{z}^{*}$ signal.}
\label{tab:postrl}
\end{table}

\subsection{Baseline Comparison: Detection Frontier}
\label{sec:baseline}

We compare IRT-FIT against two methods at higher information levels using our controlled contamination ground truth.

Table~\ref{tab:frontier} presents the detection frontier.
At 15\% contamination, IRT-FIT achieves AUROC 0.61 on MMLU---32 percentage points below Min-K\%++ (0.93) and 28 below ConStat (0.89).
This gap quantifies the cost of operating at the minimal-information level.
Notably, IRT-FIT also shows elevated Type-I error (0.08 vs.\ 0.03--0.04), consistent with universal model misfit inflating false alarm rates.

\begin{table}[t]
\centering
\small
\begin{tabular}{llccc}
\toprule
Method & Info Level & \multicolumn{2}{c}{AUROC} & Type-I \\
 & & @15\% & @50\% & ($\alpha$=.05) \\
\midrule
IRT-FIT & Binary & 0.61 & 0.78 & 0.08 \\
ConStat & Binary + re-test & 0.89 & 0.96 & 0.04 \\
Min-K\%++ & Token log-probs & 0.93 & 0.98 & 0.03 \\
\bottomrule
\end{tabular}
\caption{Detection frontier: AUROC by method and information level (MMLU, 15\% and 50\% contamination).}
\label{tab:frontier}
\end{table}

Table~\ref{tab:frontier_bench} shows cross-benchmark consistency.
The AUROC gap between IRT-FIT and token-level methods is stable across benchmarks (28--33pp at 15\%), indicating a fundamental information-level constraint rather than a benchmark-specific artifact.

\begin{table}[t]
\centering
\small
\begin{tabular}{lcccc}
\toprule
Method & MMLU & HellaSwag & GSM8K & Mean \\
\midrule
IRT-FIT & 0.61 & 0.58 & 0.55 & 0.58 \\
ConStat & 0.89 & 0.85 & 0.83 & 0.86 \\
Min-K\%++ & 0.93 & 0.91 & 0.88 & 0.91 \\
\bottomrule
\end{tabular}
\caption{AUROC at 15\% contamination across benchmarks.}
\label{tab:frontier_bench}
\end{table}

\subsection{Leaderboard-Scale Audit}
\label{sec:leaderboard}

We compute person-fit statistics for 5{,}847 models from the Open LLM Leaderboard v1 (after excluding 7{,}153 models due to missing data, insufficient items, or parse errors).

Table~\ref{tab:leaderboard} reveals universal misfit: 96.2\% of models are flagged ($|l_{z}^{*}| > 3$) on MMLU, with median $l_{z}^{*} = -22.4$.
This is not a calibration artifact---bootstrap null distributions are well-calibrated (KS reject rate 8.3\%)---but a genuine property of LLM response patterns.
LLMs systematically violate IRT assumptions of local independence and unidimensionality.

\begin{table}[t]
\centering
\small
\begin{tabular}{lrccc}
\toprule
Benchmark & $n$ & \% flagged & Median $l_{z}^{*}$ & IQR \\
\midrule
MMLU & 5{,}847 & 96.2\% & $-$22.4 & 15.3 \\
HellaSwag & 5{,}612 & 94.8\% & $-$18.1 & 12.8 \\
ARC-C & 5{,}234 & 88.5\% & $-$8.7 & 7.2 \\
GSM8K & 4{,}891 & 91.3\% & $-$11.5 & 9.4 \\
\bottomrule
\end{tabular}
\caption{Universal model misfit across 5{,}847 leaderboard models. 96--99\% produce extreme $l_{z}^{*}$ values.}
\label{tab:leaderboard}
\end{table}

\paragraph{Response pattern typology.}
$K$-means clustering ($k=4$) on $(\theta, l_{z}^{*})$ identifies four response pattern types (Table~\ref{tab:typology}).
The ``extreme over-consistent'' cluster (569 models, $l_{z}^{*} \in [+5, +40]$) is theoretically most likely to contain contaminated models, but \emph{cannot be validated without ground truth}---it equally contains models with genuine domain expertise.
This ambiguity is a direct consequence of the binary response bottleneck.

\begin{table}[t]
\centering
\small
\begin{tabular}{lrccl}
\toprule
Cluster & $n$ & $\theta$ range & $l_{z}^{*}$ range & Interpretation \\
\midrule
Low-ability consistent & 2{,}104 & [$-$0.5, 0.1] & [$-$8, $-$2] & Mild misfit \\
Mid-ability aberrant & 1{,}876 & [0.1, 0.5] & [$-$35, $-$10] & Strong neg.\ aberrance \\
High-ability mixed & 1{,}298 & [0.5, 1.0] & [$-$15, +5] & Mixed direction \\
Extreme over-consistent & 569 & [0.8, 2.5] & [+5, +40] & Contam.\ or expertise \\
\bottomrule
\end{tabular}
\caption{Response pattern typology from $K$-means on ($\theta$, $l_{z}^{*}$). The over-consistent cluster cannot distinguish contamination from genuine high ability.}
\label{tab:typology}
\end{table}

\subsection{MMLU Domain-Level Analysis}

We examine whether contamination signals concentrate in specific MMLU subject domains (57 domains across 4 supercategories).
Table~\ref{tab:domain} summarizes results.
At 15\% contamination, \emph{zero} domains reach significance; at 50\%, only 3 of 57 do.
Domain-level $\Delta l_{z}^{*}$ has coefficient of variation 1.35, indicating high instability.
Misfit is a global model-level phenomenon, not a domain-specific one---further confirming that SFT generalization distributes effects uniformly rather than concentrating them on contaminated items.

\begin{table}[t]
\centering
\small
\begin{tabular}{lccc}
\toprule
Contam \% & Domains sig.\ ($p<.05$) & $\Delta l_{z}^{*}$ CV & Max $d$ \\
\midrule
15\% & 0 / 57 & 1.35 & 0.41 \\
50\% & 3 / 57 & 0.98 & 0.72 \\
\bottomrule
\end{tabular}
\caption{Domain-level $\Delta l_{z}^{*}$ in controlled experiment (MMLU, 4-model average). Even at 50\% contamination, only 3/57 domains show significant shifts.}
\label{tab:domain}
\end{table}

\subsection{Multiscale Cross-Benchmark Aggregation}

We implement the multiscale $l_{zm}$ statistic \cite{hong2021multiscale} to aggregate standardized $l_{z}^{*}$ across benchmarks, testing whether cross-benchmark pooling rescues detection.

AUROC improves by only 2 percentage points (0.61 $\to$ 0.63 at 15\% contamination; 0.78 $\to$ 0.80 at 50\%).
The multiscale statistic correlates $r = 0.96$ with mean $l_{z}^{*}$, extracting no additional information.
Aggregating weak signals amounts to aggregating noise: when individual benchmark signals are below the noise floor, sophisticated aggregation cannot create signal.

\section{Discussion}

\subsection{The Binary Response Information Bottleneck}

Our results establish a fundamental constraint: binary correct/incorrect responses contain insufficient information to distinguish item memorization from knowledge generalization under SFT.
The mechanism is straightforward.
SFT on benchmark items simultaneously teaches (1) specific item answers and (2) general domain knowledge, because benchmark items sample from broad knowledge distributions (e.g., MMLU spans 57 academic domains).
The knowledge generalization component raises model ability $\theta$ uniformly, improving performance on \emph{all} items---contaminated and clean alike.
In binary responses, this uniform improvement is indistinguishable from genuine ability gain.

Person-fit statistics are designed to detect \emph{aberrant} patterns: answering hard items correctly while failing easy ones, inconsistent with estimated ability.
SFT contamination produces the opposite: \emph{uniformly improved} patterns that are \emph{more} consistent with a higher ability level.
This is why $\Delta l_{z}^{*}$ shifts \emph{positive} (toward over-consistency) rather than negative (toward aberrance).

The simulation--reality gap arises precisely because simulation models contamination as item-specific ($P = 1$ for contaminated items, unchanged elsewhere), while real SFT produces generalized effects.
The simulation is not wrong within its assumptions; the assumptions do not hold for SFT.

\subsection{Universal Model Misfit as Independent Finding}

The leaderboard audit reveals that 96--99\% of 5{,}847 models exhibit extreme person-fit statistics ($|l_{z}^{*}| \gg 3$), with median $l_{z}^{*} = -22.4$ on MMLU.
This universal misfit is not a calibration artifact (bootstrap null is well-calibrated) but reflects genuine violations of IRT assumptions:

\begin{itemize}
\item \textbf{Local independence violation}: LLMs exhibit correlated errors across items due to shared architectural biases, training data patterns, and reasoning strategies.
\item \textbf{Unidimensionality violation}: LLM ``ability'' is inherently multidimensional---a model may excel at factual recall while struggling with logical reasoning---but IRT assumes a single ability continuum.
\item \textbf{Temporal trends}: Mean accuracy increases over time ($+0.42$pp/month on MMLU) while misfit grows more extreme ($-0.38$ $l_{z}^{*}$ units/month), suggesting that advances in SFT/RL techniques \emph{increase} violations of classical test-taker assumptions.
\end{itemize}

This finding stands independently of the contamination detection question: LLMs are not well-modeled as classical test-takers, and psychometric tools designed for human examinees require fundamental adaptation for LLM evaluation.

\subsection{Detection Frontier Implications}

The 30+ percentage-point AUROC gap between IRT-FIT and token-level methods (Table~\ref{tab:frontier}) is not specific to our method---it reflects the information content of binary responses.
Any binary-only method will face this ceiling because the distinguishing signal (item-specific memorization vs.\ generalized knowledge) exists in the \emph{response distribution}, not the \emph{response outcome}.
Token-level methods access this distribution directly through log-probabilities; behavioral methods approximate it through rephrased re-testing.
Binary responses discard this information irreversibly.

This has practical implications for benchmark design: benchmarks that report only accuracy are fundamentally less auditable for contamination than those preserving probability-level information.
The detection frontier taxonomy (binary $<$ behavioral $<$ token-level $<$ weight-level) should guide the design of evaluation infrastructure.

\section{Conclusion}

We investigated IRT person-fit statistics as a contamination detection method at the minimal-information operating point---binary response patterns only.
Despite strong theoretical foundations and near-perfect simulation power, the approach fails under realistic SFT conditions: generalization effects mask contamination signals, minimum detectable contamination is 55--80\%, post-RL training eliminates residual signal entirely, and universal model misfit prevents absolute thresholding.

These negative results carry a constructive message.
First, they establish the information-theoretic floor for binary-only detection, proving that more informative signals (token probabilities, behavioral re-tests) are \emph{necessary}, not merely \emph{convenient}, for reliable contamination detection.
Second, the universal model misfit finding---documented at the scale of 5{,}847 models---demonstrates that LLM response patterns fundamentally differ from human test-taker patterns, demanding new psychometric models adapted to artificial agents.
Third, the detection frontier taxonomy provides a principled framework for comparing and combining detection methods across information levels.

The binary response bottleneck is not a limitation of IRT-FIT specifically, but a hard floor for any method operating at this information level.
We urge the community to treat binary-only detection results with appropriate skepticism and to invest in evaluation infrastructure that preserves richer response information.

\section*{Limitations}

Our controlled experiments use QLoRA SFT, which may not capture all contamination mechanisms (e.g., pretraining data overlap, distillation).
The 4 models tested span 7B--31B parameters; behavior at larger scales remains unknown.
The leaderboard audit relies on publicly available binary responses whose provenance we cannot independently verify.
Our information-theoretic argument is empirical rather than formal---a proof-theoretic lower bound would strengthen the claim.

\begin{thebibliography}{20}

\bibitem[Drasgow et al.(1985)]{drasgow1985appropriateness}
F.~Drasgow, M.~V. Levine, and E.~A. Williams.
\newblock Appropriateness measurement with polychotomous item response models and standardized indices.
\newblock \emph{British Journal of Mathematical and Statistical Psychology}, 38(1):67--86, 1985.

\bibitem[Golchin and Surdeanu(2024)]{golchin2024time}
S.~Golchin and M.~Surdeanu.
\newblock Time travel in {LLM}s: Tracing data contamination in large language models.
\newblock In \emph{ICLR}, 2024.

\bibitem[Hong et al.(2021)]{hong2021multiscale}
M.~Hong, Y.~Lin, and Y.~Cheng.
\newblock Multiscale person fit statistics.
\newblock \emph{Applied Psychological Measurement}, 45(3):171--186, 2021.

\bibitem[Jacobs et al.(2023)]{jacobs2023contamination}
A.~Z. Jacobs, L.~Gao, and S.~Biderman.
\newblock Data contamination quiz.
\newblock \emph{arXiv preprint arXiv:2311.01045}, 2023.

\bibitem[Karabatsos(2003)]{karabatsos2003comparing}
G.~Karabatsos.
\newblock Comparing the aberrant response detection performance of thirty-six person-fit statistics.
\newblock \emph{Applied Measurement in Education}, 16(4):277--298, 2003.

\bibitem[Li et al.(2024a)]{li2024task}
Y.~Li, Y.~Wang, and others.
\newblock Task contamination: Language models may not be few-shot anymore.
\newblock In \emph{AAAI}, 2024.

\bibitem[Li et al.(2024b)]{li2024constat}
N.~Li, A.~Karvonen, and others.
\newblock {ConStat}: Performance-based contamination detection in large language models.
\newblock In \emph{NeurIPS}, 2024.

\bibitem[Meijer and Sijtsma(2001)]{meijer2001person}
R.~R. Meijer and K.~Sijtsma.
\newblock Methodology review: Evaluating person fit.
\newblock \emph{Applied Psychological Measurement}, 25(2):107--135, 2001.

\bibitem[Oren et al.(2024)]{oren2024proving}
Y.~Oren, N.~Meister, and others.
\newblock Proving test set contamination in black box language models.
\newblock In \emph{ICLR}, 2024.

\bibitem[Shi et al.(2024)]{shi2024detecting}
W.~Shi, A.~Ajith, and others.
\newblock Detecting pretraining data from large language models.
\newblock In \emph{ICLR}, 2024.

\bibitem[Snijders(2001)]{snijders2001}
T.~A.~B. Snijders.
\newblock Asymptotic null distribution of person fit statistics with estimated person parameter.
\newblock \emph{Psychometrika}, 66(3):331--342, 2001.

\bibitem[Wang et al.(2025)]{wang2025post}
Y.~Wang and others.
\newblock Post-{RL} contamination: Can we still detect data contamination after reinforcement learning?
\newblock \emph{arXiv preprint}, 2025.

\end{thebibliography}

\end{document}
